\documentclass[12pt,a4paper]{article}

% --- Paquetes Esenciales ---
\usepackage[utf8]{inputenc}
\usepackage[T1]{fontenc}
\usepackage{amsmath}
\usepackage{amssymb}
\usepackage{graphicx}
\usepackage[spanish]{babel}
\usepackage[a4paper, margin=2.5cm]{geometry}
\usepackage{comment}

% --- Paquetes para Mejoras Visuales ---
\usepackage{hyperref}
\hypersetup{
    colorlinks=true,
    linkcolor=blue,
    filecolor=magenta,      
    urlcolor=cyan,
}
\usepackage{xcolor}
\usepackage{listings}
\usepackage{framed}

% --- Definición de Colores y Estilos ---
\definecolor{codegreen}{rgb}{0,0.6,0}
\definecolor{codegray}{rgb}{0.5,0.5,0.5}
\definecolor{codepurple}{rgb}{0.58,0,0.82}
\definecolor{backcolour}{rgb}{0.95,0.95,0.92}

\lstdefinestyle{mystyle}{
    backgroundcolor=\color{backcolour},   
    commentstyle=\color{codegreen},
    keywordstyle=\color{magenta},
    numberstyle=\tiny\color{codegray},
    stringstyle=\color{codepurple},
    basicstyle=\ttfamily\footnotesize,
    breakatwhitespace=false,         
    breaklines=true,                 
    captionpos=b,                    
    keepspaces=true,                 
    numbers=left,                    
    numbersep=5pt,                  
    showspaces=false,                
    showstringspaces=false,
    showtabs=false,                  
    tabsize=2
}
\lstset{
    style=mystyle,
    extendedchars=true,
    literate={á}{{\'a}}1 {é}{{\'e}}1 {í}{{\'i}}1 {ó}{{\'o}}1 {ú}{{\'u}}1 {Á}{{\'A}}1 {É}{{\'E}}1 {Í}{{\'I}}1 {Ó}{{\'O}}1 {Ú}{{\'U}}1 {ñ}{{\~n}}1 {Ñ}{{\~N}}1 {¿}{{?`}}1 {¡}{{!`}}1
}

\title{Guía Práctica: Implementación de un Sistema RAG (Retrieval-Augmented Generation) en Python}
\author{Prof. CERO}
\date{\today}

\begin{document}

\maketitle
\tableofcontents
\newpage

\section{Introducción al Sistema RAG}

En la era de la Inteligencia Artificial, los Modelos de Lenguaje Grande (LLMs) como GPT-4 o Llama 3 son herramientas excepcionales para el razonamiento y la generación de texto. Sin embargo, su conocimiento está limitado a la información con la que fueron entrenados y son propensos a sufrir \textit{alucinaciones} (inventar respuestas) cuando se les pregunta sobre datos privados o muy recientes.

Para solucionar esto, surge la arquitectura \textbf{RAG (Retrieval-Augmented Generation)}. Un sistema RAG actúa como un puente entre tus datos privados (apuntes, libros, repositorios) y el 'cerebro' del LLM.

\subsection{¿Cómo funciona RAG?}
Se divide en dos fases principales:
\begin{enumerate}
    \item \textbf{Retrieval (Recuperación):} Cuando el usuario hace una pregunta, el sistema busca matemáticamente en la base de datos de documentos locales y extrae únicamente los fragmentos de texto más relevantes.
    \item \textbf{Generation (Generación):} El sistema inyecta esos fragmentos en el modelo de lenguaje junto con la pregunta original. El LLM lee estos fragmentos y redacta una respuesta coherente basada \textbf{estrictamente} en el contexto proporcionado.
\end{enumerate}

\section{Implementación Paso a Paso}

A continuación, desarrollaremos un script de Python que implementa esta arquitectura paso a paso utilizando la librería \texttt{LangChain}.

\subsection{Paso 1: Carga de Documentos (Loaders)}
El primer paso es leer nuestros archivos locales (como apuntes en \texttt{.tex} o \texttt{.md}). Para ello, utilizamos un \texttt{DirectoryLoader}.

\begin{lstlisting}[language=Python]
from langchain_community.document_loaders import DirectoryLoader, TextLoader

dir_path = '/ruta/a/mis/documentos'
docs = []
for pattern in ['**/*.tex', '**/*.md']:
    loader = DirectoryLoader(
        dir_path, glob=pattern, loader_cls=TextLoader, 
        loader_kwargs={'encoding': 'utf-8'}
    )
    docs.extend(loader.load())
\end{lstlisting}
Este código escaneará todas las subcarpetas buscando archivos de texto y los cargará en la memoria RAM.

\subsection{Paso 2: Fragmentación de Texto (Splitters)}
Los modelos de lenguaje tienen un límite en la cantidad de texto que pueden leer a la vez (ventana de contexto). Por ende, debemos dividir nuestros documentos largos en pequeños párrafos superpuestos.

\begin{lstlisting}[language=Python]
from langchain_text_splitters import RecursiveCharacterTextSplitter

text_splitter = RecursiveCharacterTextSplitter(chunk_size=2000, chunk_overlap=200)
splits = text_splitter.split_documents(docs)
\end{lstlisting}
Aquí dividimos el texto en bloques de 2000 caracteres, permitiendo que 200 caracteres se superpongan entre bloques para no cortar ideas por la mitad.

\subsection{Paso 3: Vectorización y Base de Datos (ChromaDB + HuggingFace)}
Las computadoras no entienden de palabras, sino de números. Usaremos un modelo local y gratuito de \textbf{Hugging Face} para convertir nuestros fragmentos de texto en vectores matemáticos (Embeddings) y guardarlos en una base de datos vectorial llamada \textbf{ChromaDB}.

\begin{lstlisting}[language=Python]
from langchain_huggingface import HuggingFaceEmbeddings
from langchain_chroma import Chroma

embeddings = HuggingFaceEmbeddings(model_name='all-MiniLM-L6-v2')
persist_dir = 'chroma_db'

# Guardar los fragmentos como vectores en disco
vectorstore = Chroma.from_documents(
    documents=splits, 
    embedding=embeddings, 
    persist_directory=persist_dir
)
\end{lstlisting}
Esto permite que la búsqueda posterior sea instantánea.

\subsection{Paso 4: El Cerebro del Sistema (Llama 3 vía Groq)}
Para la fase de razonamiento y generación de texto, conectaremos nuestro sistema al modelo de lenguaje \textbf{Llama 3} utilizando la API ultrarrápida de \textbf{Groq}.

\begin{lstlisting}[language=Python]
from langchain_groq import ChatGroq

# Asegurate de tener configurada la variable de entorno GROQ_API_KEY
llm = ChatGroq(model='llama-3.1-8b-instant', temperature=0)
\end{lstlisting}
Establecemos la temperatura en \texttt{0} para que las respuestas sean deterministas y precisas, limitando la creatividad del modelo, lo cual es ideal para un asistente académico.

\subsection{Paso 5: Construcción de la Cadena (Chain) y el Prompt}
Finalmente, ensamblamos todo. Le diremos al sistema cómo comportarse (Prompt), le pasaremos el buscador (Retriever) y el cerebro (LLM).

\begin{lstlisting}[language=Python]
from langchain_classic.chains import create_retrieval_chain
from langchain_classic.chains.combine_documents import create_stuff_documents_chain
from langchain_core.prompts import ChatPromptTemplate

# 1. Recuperador: Buscará los 6 fragmentos más relevantes a la pregunta
retriever = vectorstore.as_retriever(search_kwargs={'k': 6})

# 2. Instrucciones para el Asistente
system_prompt = (
    'Eres un asistente académico experto. Usa los siguientes fragmentos de contexto '
    'para responder a la pregunta del usuario. Si no sabes la respuesta, di que no lo sabes.\n\n'
    '{context}'
)
prompt = ChatPromptTemplate.from_messages([
    ('system', system_prompt),
    ('human', '{input}'),
])

# 3. Ensamblar la cadena RAG
question_answer_chain = create_stuff_documents_chain(llm, prompt)
rag_chain = create_retrieval_chain(retriever, question_answer_chain)
\end{lstlisting}

\section{El Bucle Interactivo}
Para que el script actúe como un chatbot que no se cierra tras cada pregunta, implementamos un bucle \texttt{while} infinito que espera la entrada del usuario.

\begin{lstlisting}[language=Python]
print('\n¡El sistema RAG está listo! Escribe 'salir' para terminar.')
try:
    while True:
        pregunta = input('\nTu pregunta: ')
        if pregunta.lower() in ['salir', 'exit', 'quit']:
            break
            
        # Invocar la cadena y generar la respuesta
        response = rag_chain.invoke({'input': pregunta})

        print('\n--- RESPUESTA ---')
        print(response['answer'])

except (KeyboardInterrupt, EOFError):
    print('\n\nSaliendo del sistema RAG. ¡Hasta luego!')
\end{lstlisting}

¡Y listo! Con estas etapas, hemos construido un asistente académico privado, de costo cero en vectorización, y respaldado por la última tecnología en modelos de lenguaje de código abierto.

\end{document}